% === V. IDENTIFICATION METHODOLOGY ===========================================
\section{Identification Methodology}
\label{sec:method}

Figure~\ref{fig:pipeline} summarises the loop and Algorithm~\ref{alg:loop}
states it in the form needed to reimplement it; lines marked $\triangleright$
are this paper's changes to the inherited pipeline.
Sections~\ref{subsec:excitation}--\ref{subsec:iterative} carry the central
methodological contribution, the two changes to the virtual-data step.

\begin{figure*}[t]
\centering
\resizebox{0.98\textwidth}{!}{% =============================================================================
% Fig. 3 -- The corrected identification pipeline.
%
% Vector TikZ, drawn at the paper's own font sizes. Included by
% sections/method.tex inside a figure* (full text width).
%
% Layout rules, so an edit does not undo them:
%   * the four stage boxes sit on one row, equal width, gaps wide enough that
%     an edge label fits BETWEEN two boxes instead of over their borders;
%   * the three green correction boxes hang from one common baseline
%     (\fixline), so they cannot look staggered;
%   * the plant-telemetry box and the admissibility gate share a row, so the
%     orange arrow between them is a straight horizontal line;
%   * the re-nomination arrow closes the loop along the BOTTOM of the diagram.
%     Any route through the middle passes behind the gate/telemetry boxes.
%
% Colour: green = this paper's changes; blue = the gated model interface;
% orange = the plant's own force telemetry, scoring only, never in the loop.
% =============================================================================

\begin{tikzpicture}[
  font=\scriptsize,
  >={Stealth[length=1.7mm,width=1.5mm]},
  blk/.style      = {draw, rounded corners=1pt, align=flush center,
                     inner sep=3pt, text width=#1, fill=white,
                     line width=0.35pt},
  stage/.style    = {blk=30mm, fill=black!3},
  meas/.style     = {blk=26mm, fill=black!5},
  fix/.style      = {draw, rounded corners=1pt, align=flush left,
                     inner sep=2.6pt, text width=30mm, font=\tiny,
                     fill=green!12, draw=green!45!black, line width=0.4pt,
                     anchor=south},
  gate/.style     = {blk=26mm, fill=blue!10, draw=blue!60!black,
                     line width=0.7pt},
  gt/.style       = {blk=28mm, fill=orange!18, draw=orange!80!black,
                     line width=0.7pt},
  flow/.style     = {->, line width=0.5pt, black!65},
  gateflow/.style = {->, line width=0.8pt, blue!60!black},
  gtflow/.style   = {->, line width=0.8pt, dashed, orange!80!black},
  fixlink/.style  = {-, line width=0.4pt, green!45!black, dashed},
  elbl/.style     = {font=\tiny, inner sep=1pt, fill=white, midway, above}
]

\def\hgap{11mm}   % stage-to-stage gap: must clear the widest edge label

% ------------------------- measurement chain -------------------------------
\node[meas, anchor=north west] (lap) at (0,0)
  {\textbf{Recorded lap}\\[1pt]
   \SI{30}{\second} at \SI{50}{\hertz}\\
   $[v_{x},v_{y},\omega,\delta]$};
\node[meas, below=5mm of lap] (chain)
  {\textbf{Measurement chain}\\[1pt]
   achieved $\delta$; measured $m$, $l_{f}$, $l_{r}$; IMU $a_{x},a_{y}$};

% ------------------------------ stage row ----------------------------------
\node[stage, right=\hgap of lap] (s1)
  {\textbf{1. Residual learning}\\[1pt]
   $f_{NN}\!:[v_{x},v_{y},\omega,\delta]\!\mapsto\! e_{k}$\\
   corrected model $\hat f\!+\!f_{NN}$};
\node[stage, right=\hgap of s1] (s2)
  {\textbf{2. Virtual rollout}\\[1pt]
   constant $v$, steering ramp
   $\delta:0\rightarrow\delta_{\mathrm{sweep}}$};
\node[stage, right=\hgap of s2] (s3)
  {\textbf{3. Pacejka fit}\\[1pt]
   $F_{y}$ from \eqref{eq:ssforces};\\
   bounded NLLS $\rightarrow\phip$};
\node[stage, right=\hgap of s3] (s4)
  {\textbf{4. Re-nomination}\\[1pt]
   $\phip$ becomes the next\\ nominal model};

\draw[flow] (lap.east) -- ++(3mm,0) |- (s1.west);
\draw[flow] (chain.east) -- ++(3mm,0) |- (s1.west);
\draw[flow] (s1) -- node[elbl]{$\hat f\!+\!f_{NN}$} (s2);
\draw[flow] (s2) -- node[elbl]{$(\alpha,F_{y})$}    (s3);
\draw[flow] (s3) -- (s4);

% ------------------------- corrections (green) ------------------------------
% One common baseline for all three, so they cannot look staggered.
\coordinate (fixline) at ($(s1.north)+(0,8mm)$);
\node[fix] (f1) at (s1.north |- fixline)
  {\textbf{C0} friction warm start: per-axle brush $\mu$ as a \emph{floor} on
   $D_{f}$, $D_{r}$};
\node[fix] (f2) at (s2.north |- fixline)
  {\textbf{C1} rollout speed from the measured lap, floored at
   \SI{7}{\meter\per\second}; warns when \mbox{$\mu_{\mathrm{reach}}<D$}\\
   \textbf{C2} sweep bounded at $1.5\,P_{99}(|\delta_{\mathrm{obs}}|)$};
\node[fix] (f3) at (s3.north |- fixline)
  {\textbf{C3} Tikhonov target frozen at $\phip^{0}$, start point still
   tracking; multi-start, robust loss};

\draw[fixlink] (f1.south) -- (s1.north);
\draw[fixlink] (f2.south) -- (s2.north);
\draw[fixlink] (f3.south) -- (s3.north);

% ------------------- telemetry / gate row, then consumer ---------------------
% Same `below` offset for both, so the orange arrow is straight.
\node[gt,   below=11mm of s2] (truth)
  {\textbf{Plant force telemetry}\\[1pt]
   per-wheel $\alpha$, $F_{z}$, $F_{y}$ --- \emph{scoring only}};
\node[gate, below=11mm of s4] (gate)
  {\textbf{Admissibility gate}\\[1pt]
   grip ceiling vs.\ demand, understeer sign, stiffness};
\node[blk=26mm, below=6mm of gate, fill=black!3] (mpc)
  {\textbf{NMPC}\\[1pt] \texttt{acados}, $N=100$};

\draw[gateflow] (s4.south) -- (gate.north);
\draw[gateflow] (gate) -- node[elbl]{\texttt{sysid/update\_params}} (mpc);
\draw[gtflow] (truth.east) -- (gate.west);

% --------------------- re-nomination closes the loop ------------------------
\coordinate (bot) at ($(mpc.south)+(0,-6mm)$);
\draw[flow] (s4.east) -- ++(5mm,0) |- (bot) -| (s1.south);
% Sits ON the loop-back line, white-filled so it masks it -- placed at s2's x
% so it clears the elbow that rises at s1.south.
\node[font=\tiny, fill=white, inner sep=1.6pt]
  at ($(s2.south |- bot)$) {$N_{\mathrm{it}}=6$ co-identification iterations};

\end{tikzpicture}
}
\caption{The identification pipeline. Recorded lap data trains the residual
network, the corrected model is rolled out through a synthetic steering sweep,
Pacejka coefficients are fitted to that sweep, and the fit becomes the next
nominal model under relaxation. Green marks this paper's changes to the inherited loop: the friction
warm start, which measures $D$ off the recorded trace and pins it for the loop
(Section~\ref{subsec:friction}), the
excitation-aware rollout speed and the extrapolation-bounded sweep
(Sections~\ref{subsec:excitation}--\ref{subsec:iterative}), the frozen
regularisation target of \eqref{eq:fitobj}, and the relaxed coefficient update
of \eqref{eq:relax}. Blue is the gated model interface
to the controller (Section~\ref{sec:integration}); orange is the plant's own
per-wheel force telemetry, which scores the result, is armed only once the
controller has acknowledged the submitted model
(Section~\ref{subsec:deployment}), and is never admitted into the loop.}
\label{fig:pipeline}
\end{figure*}

\begin{algorithm}[t]
\footnotesize
\caption{Corrected on-track identification cycle. $\triangleright$ marks a
change to the inherited pipeline~\cite{ontrack2025}.}
\label{alg:loop}
\hrule\vspace{2pt}
\algline{in}{buffer $\mathcal{D}=\{v_{x},v_{y},\omega,\delta^{\mathrm{ach}}\}_{1:K}$ ($\SI{30}{\second}$), prior $\phip^{0}$, bounds $[\phip^{-},\phip^{+}]$}
\algline{out}{$\phip$, or \textsc{unidentified}}
\vspace{2pt}\hrule\vspace{2pt}
\algline{1}{$\delta \leftarrow$ achieved road-wheel angle, command on staleness \hfill$\triangleright$}
\algline{2}{$m,l_{f},l_{r} \leftarrow$ measured static wheel loads \hfill$\triangleright$}
\algline{3}{$\hat\mu_{f} \leftarrow \textsc{BrushFit}(\mathcal{D})$, front axle, released only if gated \hfill$\triangleright$}
\algline{4}{\textbf{if} released \textbf{then} $D^{0}_{f},D^{0}_{r} \leftarrow \hat\mu_{f}$ \quad(two-sided, both axles) \hfill$\triangleright$}
\algline{5}{$D$ held fixed for the whole loop, released or not \hfill$\triangleright$}
\algline{6}{$\phip^{\mathrm{tik}} \leftarrow \phip^{0}$ \quad(frozen for all iterations) \hfill$\triangleright$}
\algline{7}{$v \leftarrow \mathrm{clip}(\overline{v_{x}}(\mathcal{D}),\,v_{\min},\,\infty)$ \hfill$\triangleright$}
\algline{8}{$\delta_{\mathrm{sweep}} \leftarrow \mathrm{clip}(1.5\,P_{99}(|\delta|),\,\delta_{\min},\,\delta_{\max})$ \hfill$\triangleright$}
\algline{9}{$\mu_{\mathrm{reach}} \leftarrow v^{2}\delta_{\mathrm{sweep}}/(gL)$; \textbf{warn if} $\mu_{\mathrm{reach}}<\max_{a} D^{0}_{a}$ \hfill$\triangleright$}
\algline{10}{$\phip \leftarrow \phip^{0}$}
\algline{11}{\textbf{for} $i=1$ \textbf{to} $N_{\mathrm{it}}$ \textbf{do}}
\algline{12}{\algin{1.5}$f_{NN} \leftarrow \textsc{TrainResidual}(\mathcal{D},\phip)$ \quad(re-initialised each $i$)}
\algline{13}{\algin{1.5}$\mathcal{S} \leftarrow \textsc{Rollout}(\phip,f_{NN},v,\delta_{\mathrm{sweep}})$}
\algline{14}{\algin{1.5}$\tilde\phip \leftarrow \arg\min_{B,C,E} \|F_{y}(\phip;\mathcal{S})-\hat F_{y}(\mathcal{S})\|^{2} + \lambda\|\phip-\phip^{\mathrm{tik}}\|^{2}$, $D$ held \hfill$\triangleright$}
\algline{15}{\algin{1.5}\quad subject to $\phip\in[\phip^{-},\phip^{+}]$, started from $\phip$ of $i-1$}
\algline{16}{\algin{1.5}$\phip \leftarrow (1-\beta)\,\phip + \beta\,\tilde\phip$ \quad($\beta=\num{0.20}$) \hfill$\triangleright$}
\algline{17}{\textbf{if} any $\phip$ component on a bound \textbf{then} report the degeneracy \hfill$\triangleright$}
\algline{18}{\textbf{if not} \textsc{Gates}$(\phip)$ \textbf{then return} \textsc{unidentified} \hfill$\triangleright$}
\algline{19}{submit $\phip$; keep the previous model until acknowledged \hfill$\triangleright$}
\vspace{2pt}\hrule\vspace{2pt}
\algline{}{\textsc{Gates}: axle stiffness sign and magnitude; $l_{f}C_{f}<l_{r}C_{r}$ and $v_{\mathrm{crit}}$ outside the speed range; grip ceiling $\ge$ the trajectory's lateral demand (Section~\ref{sec:integration}).}
\end{algorithm}

\subsection{Residual model}
\label{subsec:residual}

The residual carries whatever the nominal model
\eqref{eq:vydot}--\eqref{eq:mf} does not represent. How that map is
parameterised is a modelling decision: known vehicle physics shapes its input
space and, optionally, its objective (Section~\ref{subsec:physinputs}), so that
the network's small capacity is spent on what is genuinely unmodelled rather
than on rediscovering structure \eqref{eq:slip} and
\eqref{eq:vydot}--\eqref{eq:omdot} already state exactly.

The default residual is the baseline multilayer perceptron: 4 inputs, one
hidden layer of 8 units, 2 outputs, LeakyReLU, 58 parameters, full-batch
Adam~\cite{kingma2015adam} at $5\times10^{-4}$ on an MSE objective, in
PyTorch~\cite{paszke2019pytorch}. Because the full-scale residual has different
structure from the scaled platform's, the architecture is selectable:
\texttt{baseline}, \texttt{wide} (16 hidden units, 114 parameters, weight
decay), \texttt{physics\_inputs} (74 parameters), \texttt{ensemble} (three
averaged members, 174 parameters), and \texttt{s4} (\numrange{102}{110}
parameters). The results below use the baseline MSE objective, so that the
identifiability statements of
Sections~\ref{subsec:excitation}--\ref{subsec:iterative} are not confounded
with a change of residual.

\subsubsection{Physics-augmented inputs and objective}
\label{subsec:physinputs}
The \emph{input space} is augmented with the axle slip angles of
\eqref{eq:slip}: since the tire law \eqref{eq:mf} is a function of
$(\alpha,F_{z})$ and of nothing else in $\vecz$, supplying
$\alpha_{f},\alpha_{r}$ directly spares the network a relation the model
already states. The \emph{objective}
\begin{equation}
\mathcal{L}=\mathcal{L}_{\mathrm{data}}
+\lambda_{s}\mathcal{L}_{\mathrm{steady}}
+\lambda_{m}\mathcal{L}_{\mathrm{sym}}
+\lambda_{g}\mathcal{L}_{\mathrm{smooth}},
\label{eq:pinnloss}
\end{equation}
with $\lambda_{s}=0.1$, $\lambda_{m}=0.05$, $\lambda_{g}=10^{-4}$, adds three
penalties to the one-step MSE~\cite{raissi2019pinn,deepdynamics2024}, each
answering a weakness of training on one lap of ordinary driving.
$\mathcal{L}_{\mathrm{steady}}$ keeps the two residual channels jointly
consistent with \eqref{eq:vydot}--\eqref{eq:omdot} on a quasi-steady mask, so
they cannot imply an axle force pair no tire could produce --- and it is that
force pair, not the state prediction, that stage~3 fits.
$\mathcal{L}_{\mathrm{sym}}$ enforces the odd symmetry $e(-\vecz)=-e(\vecz)$ of
a laterally symmetric vehicle, supplying from geometry the half of the input
domain a circuit under-covers. $\mathcal{L}_{\mathrm{smooth}}$ is a soft
Lipschitz bound on how fast the residual may grow inside the range the
extrapolation bound of Section~\ref{subsec:iterative} allows. Derivations,
masks and the precomputation that makes the nominal-model half of
\eqref{eq:pinnloss} free are in the supplementary material.

One property of \eqref{eq:pinnloss} is not a hyperparameter.
$\mathcal{L}_{\mathrm{data}}$ is a mean square over a one-step increment in
\si{\meter\per\second} and \si{\radian\per\second}, whereas the two balance
residuals of \eqref{eq:vydot}--\eqref{eq:omdot} are a force and a moment; in
those units the terms are four to six orders of magnitude apart and
$\lambda_{s}$ is not a relative weight at all. Measured at initialisation on a
\SI{30}{\second} buffer of this vehicle, the unnormalised steady term
contributes \SI{99.9996}{\percent} of $\mathcal{L}$ against
\SI{0.0004}{\percent} for the data term, a ratio of $\num{2.3e5}$: such a
residual is fitted to cancel the nominal model's own force imbalance and never
to the data it was given. Each residual is therefore expressed as the state
increment it implies, the unit the network's output is in,
\begin{equation}
r_{\mathrm{lat}}=\frac{T_{s}}{m}\,\Big(m\dot v_{y}+mv_{x}\omega-(\Fyr+\Fyf\cos\delta)\Big),
\label{eq:physnorm}
\end{equation}
and $r_{\mathrm{yaw}}$ likewise with $T_{s}/I_{z}$, before
$\mathcal{L}_{\mathrm{steady}}=\mathrm{mean}(r_{\mathrm{lat}}^{2})+\mathrm{mean}(r_{\mathrm{yaw}}^{2})$
is formed. The shares then become \SI{98.53}{\percent} data and
\SI{1.47}{\percent} steady, recovering the meaning \eqref{eq:pinnloss} gives
$\lambda_{s}$. This is a correction to the objective as published, not an
implementation detail.

\subsubsection{Temporal residual}
\label{subsec:s4}
The \texttt{s4} architecture replaces the memoryless MLP with an S4D diagonal
state-space residual~\cite{gu2022s4d} over a sliding window of $W=20$ samples
(\SI{0.7}{\second} at $T_{s}=\SI{0.035}{\second}$), so the residual can carry
tire-relaxation and yaw-damping dynamics a per-sample map cannot. We claim
neither the architecture, which is also in concurrent
work~\cite{visionsysid2026}, nor the \num{8.7}$\times$ reduction in its
training cost, an exact algebraic rewrite of the same forward pass that changes
no identified coefficient; both are derived and measured in the supplementary
material. Its accuracy effect is positive but measured offline
(Section~\ref{subsec:ablations}), and the identifiability result of
Section~\ref{subsec:excitation} holds for either architecture.

\subsection{Excitation-aware virtual rollout}
\label{subsec:excitation}

This is the central correction. Consider stage~2 of
Section~\ref{subsec:baseline}: a constant-speed rollout at $v$ with the
steering ramped to $\delta_{\mathrm{sweep}}$. In quasi-steady cornering the
lateral acceleration a single-track vehicle develops kinematically is
$a_{y}=v^{2}\kappa$ with $\kappa\approx\delta/L$, so the largest friction
coefficient the rollout can \emph{demand} of the tire is
\begin{equation}
\boxed{\;\mu_{\mathrm{reach}}\;\approx\;\frac{v^{2}\,\delta_{\mathrm{sweep}}}{g\,L}\;}
\label{eq:mureach}
\end{equation}
Equation~\eqref{eq:mureach} contains no tire parameter: it is a property of the
rollout configuration alone, and an \emph{upper} bound rather than an equality,
because $\kappa\approx\delta/L$ is the Ackermann curvature while the
steady-state curvature a compliant vehicle holds,
$\kappa=\delta/(L+K_{us}v^{2}/g)$, is smaller. Reading it as a ceiling is all
the argument needs: if even the ceiling sits below $D$, the peak is certainly
not reached, the synthetic data never leaves the linear ramp of \eqref{eq:mf},
only the product $BCD$ is identifiable
(Section~\ref{subsec:identifiability}), and the bounded optimiser resolves the
residual freedom on a box edge --- the signature measured in
Section~\ref{subsec:resultstire}.

The baseline implementation selects the rollout speed as
$v=\mathrm{clip}(\overline{v_{x}},\,2,\,4)\,\si{\meter\per\second}$,
representative of a 1:10 platform. On our vehicle, with
$\delta_{\mathrm{sweep}}=\SI{0.4}{\radian}$ and $L=\SI{1.5334}{\meter}$,
\eqref{eq:mureach} gives $\mu_{\mathrm{reach}}=0.43$ at
\SI{4}{\meter\per\second}: the rollout cannot ask the tire for more than
\SI{0.43}{\g} whatever the tire is. At 1:10 scale, with
$L\approx\SI{0.32}{\meter}$, the same speed yields
$\mu_{\mathrm{reach}}\approx2.0$ and the pathology is invisible.

\subsubsection{Correction}
\label{subsec:rollout}
The rollout speed is now taken from the speed the vehicle was actually driven
at, clamped by configuration bounds (lower bound \SI{15}{\meter\per\second},
no upper bound), and the pipeline warns whenever \eqref{eq:mureach} evaluates
below the prior's $D$, turning a silent degeneracy into a startup diagnostic.
The same bound sets a data-collection speed: identification data should be
gathered at $v\ge\sqrt{gLD/\delta_{\mathrm{sweep}}}$.

\subsection{Bounded extrapolation and frozen-prior regularisation}
\label{subsec:iterative}

Raising the rollout speed exposes a second, opposite failure: the rollout
queries the residual network at steering angles it has never seen, the
observed $|\delta|$ on a pure-pursuit lap staying below
$\approx\SI{0.06}{\radian}$ while the sweep ran to \SI{0.4}{\radian}. The fit
reads the network's extrapolation as additional grip, and because each
iteration's answer became both the next iteration's nominal model \emph{and}
its regularisation target, the error compounds: $D$ walked from the
\num{1.009}/\num{1.002} prior to \num{1.5402}/\num{1.8721} over six iterations,
a claimed \SI{1.7}{\g} on a vehicle measuring \SI{1.0}{\g} in steady state.

Two changes close the loop:
\begin{enumerate}
\item \textbf{Extrapolation bound.} The sweep terminates at
$\delta_{\mathrm{sweep}}=\mathrm{clip}\big(\lambda_{e}\cdot
P_{99}(|\delta_{\mathrm{obs}}|),\,\delta_{\min},\,\delta_{\max}\big)$ with
$\lambda_{e}=1.5$, $\delta_{\min}=\SI{0.05}{\radian}$,
$\delta_{\max}=\SI{0.54}{\radian}$, so the network is never queried more than
\SI{50}{\percent} beyond the steering range it was trained on. Both
$\delta_{\max}$ and the inherited $\delta_{\mathrm{sweep}}=\SI{0.4}{\radian}$
lie above the plant's \SI{16.0}{\degree} ($\SI{0.279}{\radian}$) mechanical
lock, by a factor of \num{1.9} and by \SI{43}{\percent} respectively. That is
deliberate: the sweep is a \emph{virtual} rollout, and clamping it to the lock
would cap $\mu_{\mathrm{reach}}$ a second time, by the actuator instead of the
rollout speed. In practice the binding limit is the data-driven
$\lambda_{e}P_{99}(|\delta_{\mathrm{obs}}|)$ term. The residual is nonetheless
evaluated outside the plant's reachable steering set, the extrapolation
recorded in Section~\ref{subsec:assumptions}, A6.
\item \textbf{Frozen prior.} The Tikhonov target is frozen at the configured
prior $\phip^{0}$ for all iterations, while the optimiser's \emph{start point}
continues to track the previous iteration. The per-axle objective is
\begin{equation}
\min_{\phi}\;\sum_{k}\rho\!\big(r_{k}(\phi)\big)\;+\;\sum_{j}w_{j}\big(\phi_{j}-\phi^{0}_{j}\big)^{2},
\label{eq:fitobj}
\end{equation}
with $r_{k}$ the force residual, $\rho$ a soft-$\ell_{1}$ robust loss and
$w_{j}$ scaled to the data residual budget by a single dimensionless weight
$\lambda=0.05$. The regulariser exists to hold the coefficients the manoeuvre
does not excite, $E$ above all, at their nominal value instead of on a box
edge.
\end{enumerate}

\subsection{What the rollout fit identifies, and why the loop must be damped}
\label{subsec:contraction}

Two further properties belong to the loop itself, and together they account for
everything the corrected pipeline does with $D$ and with the linear range.

\subsubsection{The rollout fit sees only the product $BCD$}
\label{subsec:bcd}
In the virtual-data step the low-slip degeneracy of
Section~\ref{subsec:identifiability} is stronger, and structural rather than
excitation-dependent. Stage~3 does not observe a force; it reconstructs one
from the rollout's own quasi-steady yaw balance \eqref{eq:ssforces}, which on
the synthetic trajectory reduces to
\begin{equation}
\frac{F_{y,f}}{F_{z,f}}=\frac{v_{x}\omega}{g\cos\delta},\qquad
\frac{F_{y,r}}{F_{z,r}}=\frac{v_{x}\omega}{g},
\label{eq:rollobs}
\end{equation}
so the only quantity \eqref{eq:fitobj} ever compares against is the lateral
acceleration the rollout realises --- produced by the \emph{previous}
coefficient set through its slope $F_{z}BCD$. The fit's residual therefore
depends on $(B,C,D)$ only through their product, and $D$ is the free coordinate
of the hyperbola $BCD=\mathrm{const}$: not separable from $B$, and stopping
wherever the box constraint the other factors reach first puts it. Replayed on
five identification buffers of a live $\mu=\num{1.05}$ run
(Section~\ref{subsec:ablations}), the product is recovered tightly and
repeatably, $BCDF_{z,f}$ over \SIrange{18.9}{21.1}{\kilo\newton\per\radian}
against a plant measured at \SI{20.4}{\kilo\newton\per\radian} from a prior
\SI{35}{\percent} below it, while the $D$ reaching that same product ranges
over \numrange{0.69}{2.00}.

$D$ must therefore not be asked of this fit at all. It is measured on the
\emph{recorded} trace by the brush estimator of
Section~\ref{subsec:friction}, which fits $F_{y}(\alpha)$ directly and is not
subject to \eqref{eq:rollobs}, then held fixed while $B$, $C$ and $E$ carry the
stiffness the rollout does determine (Algorithm~\ref{alg:loop}, lines 3--5
and 14).

\subsubsection{The loop is not a contraction}
\label{subsec:relax}
Stage~4 makes each iteration's answer the nominal model stage~2 rolls out, and
stage~3 then fits the trajectory that model produced, over a steering range
already extrapolating $\lambda_{e}=1.5\times$ past the observed one. The fit
reads its own forces back, so the loop is a positive feedback path on its own
extrapolation error rather than a contraction, and iterating it is not a
refinement: replayed offline on five recorded buffers, axle force RMSE against
the plant's telemetry degrades monotonically with iteration count, \num{38.9},
\num{53.0}, \num{58.5} and \SI{70.4}{\newton} at the front for
$N_{\mathrm{it}}\in\{1,2,3,6\}$. Pinning $D$ does not remove this; it relocates
it into $B$, $C$ and $E$.

The remedy is to damp the step, not to shorten the loop. Stage~4 becomes a
relaxed fixed-point iteration,
\begin{equation}
\phip^{(i)} \leftarrow (1-\beta)\,\phip^{(i-1)} + \beta\,\tilde\phip^{(i)},
\label{eq:relax}
\end{equation}
with $\tilde\phip^{(i)}$ the solution of \eqref{eq:fitobj} at iteration $i$ and
$\beta=\num{0.20}$; $\beta=1$ recovers the inherited update exactly. The sweep
over $\beta$ (Section~\ref{subsec:ablations}) carries the control the
contribution needs: $\beta=0$ freezes the coefficients at the pinned prior,
performing no identification at all, and is the worst setting. The iterations
do carry information; they have to be damped for it to survive.

\subsection{Pacejka fit}
\label{subsec:fit}

Fitting applies the steady-state force recovery of \eqref{eq:ssforces} to the
synthetic data under box constraints
$\phip\in([4.0,1.2,0.4,-3.0],[20.0,2.2,2.0,1.0])$ per axle, over $B$, $C$ and
$E$ with $D$ held at the value Section~\ref{subsec:friction} put there.
Relative to the baseline we (i) raised the jittered multi-starts from 1 to 8
and (ii) added a derivative-free simplex~\cite{neldermead1965} and a global
differential-evolution solver~\cite{storn1997de} alongside trust-region
reflective~\cite{branch1999trust}, all through the same optimisation
library~\cite{virtanen2020scipy}; population-based global search on
Magic-Formula coefficients follows Cabrera \emph{et
al.}~\cite{cabrera2004genetic}. The prior carries the measured tire of
Table~\ref{tab:ref}. One caveat on (i): \texttt{num\_starts} is honoured by the
trust-region and simplex solvers only, differential evolution being a
population search that ignores it, so the multi-start claim describes the
\texttt{Baseline-NN-MSE} arm and not the shipped \texttt{Ours} arm. Neither
change addresses the degeneracy, and Section~\ref{subsec:sweep} verifies both
as secondary: at a \SI{4}{\meter\per\second} rollout every tire tested fits
back as $\hat D\approx\num{0.5}$ with coefficients railed, and the answer moves
neither with the multi-start seed nor with a jitter seven times wider.

\subsection{Measurement-side corrections}
\label{subsec:measurement}

Three corrections apply to the measurement chain rather than to the estimator.

\subsubsection{Achieved versus commanded steering}
$\alpha_{f}$ in \eqref{eq:slip} was built from the \texttt{/drive} setpoint.
The measured achieved/commanded ratio is \num{0.76}, so $\alpha_{f}$ was
systematically $\approx\SI{25}{\percent}$ too large and the fitted front
stiffness $\approx\SI{20}{\percent}$ too small
(\SI{13725}{\newton\per\radian} against \SI{17198}{\newton\per\radian} fitted
to the simulator's own per-wheel forces over the same run); using
\texttt{feedback/\allowbreak steering\_angle} recovers
\SI{18486}{\newton\per\radian}. All three are effective slopes over that run's
load range and so fall below the $\bar F_{z}=1$ value of Table~\ref{tab:ref};
the comparison is between them.

\subsubsection{Mass and geometry}
$m$ and $l_{wb}$ are taken from the measured static wheel loads
(Table~\ref{tab:vehicle}) rather than from the configuration file, and
$l_{wb}\equiv l_{f}+l_{r}$ is enforced by a regression test, because
\eqref{eq:ssforces} uses the wheelbase for the normal loads and the axle split
for the force division, and the two disagreeing is a silent bias.

\subsubsection{Proprioceptive friction warm-start}
\label{subsec:friction}
The inherited pipeline had no prior for the peak factor: $D$ started from a
configuration constant and was left to the fit. What follows recovers one from
the sensors the vehicle already carries, in four steps.

\paragraph{Axle forces and slip angles without a force sensor}
Newton--Euler on the single-track body of
\eqref{eq:vydot}--\eqref{eq:omdot}, solved for the axle forces rather than the
accelerations, gives
\begin{equation}
F_{y,f}=\frac{I_{z}\dot\omega+l_{r}m a_{y}}{L\cos\delta},\qquad
F_{y,r}=\frac{l_{f}m a_{y}-I_{z}\dot\omega}{L},
\label{eq:newtoneuler}
\end{equation}
which needs the IMU's lateral acceleration and the yaw acceleration and no tire
model. Vertical loads carry the longitudinal transfer
$F_{z,f/r}=m(g\,l_{r/f}\mp a_{x}h_{cg})/L$ \cite{rajamani2012}. Yaw
acceleration is the derivative of a local polynomial fit to the yaw rate
(Savitzky--Golay, window 21, order 2) \cite{savitzky1964} rather than a
difference, because $I_{z}\dot\omega$ is a $\approx\SI{10}{\percent}$ term in
$F_{y,f}$ here and a differenced odometry yaw rate puts its noise straight into
it. Slip angles come from \eqref{eq:slip} on the buffer already collected.
Without an IMU the accelerations are differentiated from the states by the same
filter; with one, the RMS difference between the two is logged as a bias
diagnostic.

\paragraph{Joint fit of stiffness and peak friction}
Per axle, cornering stiffness and peak friction are fitted jointly to the Fiala
brush model~\cite{fiala1954,svendenius2004brush,pacejka2012book}
\begin{equation}
F_{y}=\mu F_{z}\big(1-(1-\xi)^{3}\big)\mathrm{sgn}\,\alpha,\quad
\xi=\min\!\Big(\frac{C_{\alpha}|\tan\alpha|}{3\mu F_{z}},1\Big),
\label{eq:brush}
\end{equation}
by bounded nonlinear least squares with a soft-$\ell_{1}$ loss
\cite{virtanen2020scipy,branch1999trust}. The brush model replaces
\eqref{eq:mf} here because it has two parameters, both of which the manoeuvre
can support, against the four of the model whose non-identifiability at low
slip this paper is about. A two-stage fit is avoided because estimating
$C_{\alpha}$ first, on data already slightly nonlinear, biases it.

What beats the utilisation bound is that $\mu$ enters \eqref{eq:brush} through
the \emph{curvature} of $F_{y}(\alpha)$ and not only its asymptote. Writing
$r=|F_{y}|/(C_{\alpha}|\tan\alpha|)$ for the force's departure from its own
linear extrapolation, \eqref{eq:brush} inverts to
\begin{equation}
\mu=\frac{2k}{3-\sqrt{3(4r-1)}},\qquad k=\frac{C_{\alpha}|\tan\alpha|}{3F_{z}},
\label{eq:brushinv}
\end{equation}
so a \SI{30}{\percent} departure from linearity ($r=0.7$) already pins $\mu$,
well below the limit. Equation~\eqref{eq:brushinv} is used only for the
excitation diagnostics.

\paragraph{When the estimate is released}
Equation~\eqref{eq:brushinv} states its own failure mode:
$\mathrm{d}\mu/\mathrm{d}r$ grows without bound as $r\to1$, i.e.\ as the curve
stays linear. The estimate is released only when four conditions hold:
absolute grip demand $\max|F_{y}|/F_{z}\ge\num{0.40}$, peak friction
utilisation $\max|F_{y}|/(\hat\mu F_{z})\ge\num{0.40}$, relative precision
$\sigma_{\hat\mu}/\hat\mu\le\num{0.15}$ from the Gauss--Newton covariance, and
$\hat\mu$ not resting on its own search bound; otherwise the module logs which
condition failed and falls back to the utilisation bound.

The first two are the same quantity with and without the fit in the
denominator, and both are needed. The utilisation test divides by the
\emph{fitted} $\hat\mu$ of the same least-squares call, so it is circular: an
axle whose $\mu$ comes out too low reports a correspondingly high utilisation
and clears the gate on its own error --- on a live run a buffer realising only
\num{0.31} of $F_{z}$ returned $\hat\mu=\num{0.539}$ at utilisation \num{0.58}
and passed, against a plant of \num{1.05}. The absolute ratio has no such
freedom, and on the five buffers of that run it separates usable from unusable
exactly: $\{\num{0.146},\num{0.314}\}$ rejected,
$\{\numrange{0.541}{0.677}\}$ accepted. The precision and railing conditions
remove what survives both. The \SI{40}{\percent} floor is not tuned here; it is
the level the lateral-dynamics friction-estimation literature converges on, as
surveyed by Acosta \emph{et al.}~\cite{acosta2017friction}, who report
estimators failing below $|a_{y}|=\SI{0.4}{\g}$ ``due to the proximity of the
tyre force curves in the low slip regions'' and a \SIrange{40}{80}{\percent}
utilisation working range. Samples with $|a_{x}|>\SI{0.35}{\g}$ are dropped
rather than corrected, since the axle longitudinal force split a
friction-ellipse correction needs is not observable from this sensor set;
samples below \SI{2.5}{\meter\per\second} are dropped because \eqref{eq:slip}
is ill-conditioned there; and samples whose force opposes their own slip angle
are dropped as sign errors or noise.

\paragraph{Applied two-sided, then pinned}
\label{subsec:frictionfloor}
The released $\hat\mu$ \emph{replaces} $D$ on both axles, in either direction,
and $D$ is then held fixed for the whole co-identification loop. Four points
fix that rule.

First, both axles take the \emph{front} fit. The rear slip angle does not carry
$\delta$, so on this trajectory the rear force curve has too little curvature
to support \eqref{eq:brushinv} and its $\mu$ comes out high:
\num{1.495}/\num{1.383}/\num{1.257} rear against
\num{1.011}/\num{0.964}/\num{0.837} front on the three admitted buffers of a
$\mu=\num{1.05}$ run, with $\sigma_{\hat\mu}/\hat\mu=\num{0.105}$ at the rear,
inside the gate. This is not a yaw-inertia error: in quasi-steady cornering
$F_{y,f}/F_{y,r}\to l_{r}/l_{f}$, the same ratio as $F_{z,f}/F_{z,r}$, so
$\mu_{f}=\mu_{r}$ whatever $I_{z}$ is, and sweeping $I_{z}$ over
\SIrange{30}{80}{\kilo\gram\meter\squared} only widens the axle gap from
\num{0.368} to \num{0.550}. There is one road surface, so the axle the
manoeuvre identifies speaks for both.

Second, it is two-sided rather than a floor: the floor rule it replaces was
protecting against a $D$ that \emph{moved} between cycles, not one that was
low, and pinning supplies that directly without a one-way bias a low-friction
surface makes untenable. Third, it is pinned and not merely seeded, since
\eqref{eq:rollobs} leaves $D$ free along $BCD=\mathrm{const}$ and on a live run
a \num{0.785} seed ended at the \num{2.000} bound within one cycle. Fourth, an
axle with no released $\hat\mu$ holds the $D$ it came in with rather than
reverting to fitting it, since Section~\ref{subsec:bcd} shows there is nothing
in the rollout to fit it against.

Two downstream mechanisms make the fixed $D$ what the interface needs. The
frozen Tikhonov target of \eqref{eq:fitobj} is taken from the warm-started
prior, so whatever lands there is both initial guess and regularisation target
for all $N_{\mathrm{it}}$ iterations; and the controller re-ingests the
reference speed profile at
$\min(a_{y,\max}^{\mathrm{cfg}},\,a_{y}^{\max}\eta)$ on every accepted update,
so a $D$ moving \emph{within} a cycle rewrites the reference under the running
controller. The requirement is that $D$ not move during a cycle, not that it
only ever rise.

Pinning puts the whole weight of the result on the gate, and
$\sigma_{\mu}/\mu$ from the fit covariance is a \emph{precision} figure, not an
accuracy one: a wrong mass, wrong yaw inertia, IMU bias or road bank all fit
tightly and wrongly through \eqref{eq:newtoneuler}, axle force scaling linearly
with $m$. Here the configuration file's \SI{240.0}{\kilo\gram} against the
\SI{269.6}{\kilo\gram} the wheels actually carry biased $\mu$
\SI{11}{\percent} low until the measured value replaced it
(Table~\ref{tab:vehicle}); likewise the configured wheel coefficient is not the
effective one the physics step uses, the server multiplying it by the road
surface's factor of \num{0.70}, so an estimate validated against the former
would be judged \SI{30}{\percent} short while being correct. Of the inputs to
\eqref{eq:newtoneuler} only $I_{z}$ now remains unmeasured; it enters
$F_{y,f}$ at the \SI{10}{\percent} level (Section~\ref{subsec:assumptions}).

\subsection{Deployment}
\label{subsec:deployment}

The identifier runs as a single ROS~2 node whose acquisition and state-machine
timer ticks at \SI{50}{\hertz} (the \SI{20}{\hertz} quoted elsewhere is the
controller's period), cycling a rolling FIFO buffer through four states:
collect, train, await the controller's acknowledgement, run, re-identifying
every $T_{\mathrm{reid}}=\SI{30}{\second}$ from the newest buffer. Three
interface details matter for reproduction. The achieved road-wheel angle is
used where published, falling back to the commanded setpoint on staleness with
the substitution logged. The friction warm start runs on \emph{every} cycle:
restricting it to the first, as an earlier version did, does not hold $D$
still, since the node re-reads the static \texttt{pacejka\_model} prior each
cycle, so every later cycle cold-starts $D$ from a value unrelated to the road
--- on a live run cycles 2--6 then walked $D$ to its \num{2.0} bound. And the
coefficients are submitted over a service call whose non-acknowledgement leaves
the previous model in place, so a rejected identification degrades to the last
accepted model rather than to none; the peak-friction outputs are published on
a latched diagnostic topic so a late-starting consumer sees the current value.

That acknowledgement also arms the scoring: the node's own online metrics and
the parameter forward to the passive benchmark node both fire from the
acknowledgement handler, not from the end of training. A submission the
controller \emph{rejects} on its admissibility gate
(Section~\ref{sec:integration}) is a model the car never drove on, and the
seconds between a fit completing and its acknowledgement landing are seconds
still driven on the previous model; scoring either attributes to one model the
behaviour of another.

\subsection{Validation protocol}
\label{subsec:validation}

Three levels of validation are used, in increasing strength:
\begin{enumerate}
\item \textbf{One-step prediction} on a held-out run, as in the baseline: RMSE
and $R^{2}$ of $v_{y}$ and $\omega$ under
$\hat\vecx_{k+1}=f(\vecx_{k},\vecu_{k},\phip)$. A dedicated ROS~2 node extends
this to configurable multi-step horizons (1/5/10 steps) with Euler/Heun/RK4
integration and online statistics in Welford's form~\cite{welford1962}.
Samples are admitted only above \SI{1}{\meter\per\second}, the gate the
identification data uses. This is not a division-by-zero guard: since
$\alpha=\arctan(v_{y}/v_{x})$, a few tenths of a metre per second of lateral
drift saturates the tire curve and the model predicts yaw rates of several
\si{\radian\per\second} against a ground truth of a few tenths, so a looser
gate lets the low-speed bootstrap dominate a metric meant for the racing
regime.
\item \textbf{Direct force comparison} against
\texttt{feedback/\allowbreak tire\_forces}: the identified curve is evaluated at
the measured $(\alpha,F_{z})$ and compared per axle against the measured
$F_{y}$, and against the plant's own curve~\eqref{eq:physx} at the same points
(Table~\ref{tab:ref}), which separates the error of the fit from the error of
representing a PhysX tire.
\item \textbf{Closed-loop tracking} under the MPC using the identified model:
mean and maximum lateral deviation, its time-weighted integral, heading error
and solver cost, split by active controller so the pure-pursuit acquisition
phase is not scored against a model it does not use. Lap counting runs from the
first crossing of the raceline's $s=0$, so the drive from the spawn point is an
untimed out-lap.
\end{enumerate}

All three levels score the \emph{accepted} model only, for the reason given in
Section~\ref{subsec:deployment}.
